\documentclass[12pt,a4paper]{article}
        \makeatletter
        \def\input@path{{../}}
        \makeatother
        \usepackage{almeria-fichas}

        \FichaMeta{Sesion 01}{¿Que es una guia turistica matematica?}{Bloque 1 · Mapa y coordenadas}{Comprender el reto del proyecto y descubrir que una ciudad tambien puede describirse con datos, medidas y preguntas matematicas.}{Conceptual}{Recordar, Comprender, Analizar}

        \begin{document}

        \FichaCabecera

        \begin{materialesbox}
        \begin{itemize}
    \item Cuaderno o ficha impresa
    \item Lapiz y goma
    \item Lapis de colores
    \item Plano o imagen de Almeria si el profesorado la facilita
\end{itemize}
        \end{materialesbox}

        \begin{pasosbox}
        \begin{enumerate}
    \item Lee el titulo y el objetivo en voz baja.
    \item Subraya las palabras clave: guia, ciudad, dato, medida.
    \item Resuelve primero los ejercicios de recordar y comprender.
    \item Deja para el final los ejercicios de analizar.
\end{enumerate}
        \end{pasosbox}

        \begin{recordatoriobox}
        \begin{itemize}
    \item Una guia turistica informa y orienta.
    \item Las matematicas permiten medir, comparar y representar.
    \item Un buen dato lleva numero y, si hace falta, unidad.
\end{itemize}
        \end{recordatoriobox}

        \begin{apoyobox}
        \begin{itemize}
    \item Si una consigna te parece larga, divide la lectura en dos partes.
    \item Puedes hacer un esquema antes de escribir una respuesta completa.
    \item Si usas un dato numerico, explica tambien que significa.
\end{itemize}
        \end{apoyobox}

        \BloomSection{recordar}

\BloomExercise{recordar}{1}{Elementos de una guia}{Escribe \textbf{cuatro elementos} que sueles encontrar en una guia turistica. Por ejemplo: mapa, horario o precio.}{5}

\BloomExercise{recordar}{2}{Datos de la ciudad}{Anota \textbf{tres lugares} de Almeria y escribe al lado \textbf{un dato matematizable} de cada uno. Ejemplo: altura, distancia, numero de visitantes.}{6}

\BloomSection{comprender}

\BloomExercise{comprender}{1}{Descriptivo o cuantitativo}{Lee estos datos y clasificalos: \emph{La Alcazaba es antigua}; \emph{La playa mide 900 m}; \emph{La catedral tiene una torre alta}; \emph{El paseo dura 15 minutos}. Explica por que.}{6}

\BloomExercise{comprender}{2}{Para que sirve el numero}{Escribe dos frases: una con un dato \textbf{solo descriptivo} y otra con un dato \textbf{numerico}. Despues explica cual ayuda mas a planificar una visita y por que.}{6}

\BloomSection{analizar}

\BloomExercise{analizar}{1}{Preguntas matematicas}{Elige \textbf{dos lugares} de Almeria. Para cada uno, redacta \textbf{dos preguntas matematicas} que podrian aparecer en una guia. No hace falta resolverlas.}{7}

\BloomExercise{analizar}{2}{Seleccion de informacion}{Un equipo quiere poner en la guia estos datos: color de una puerta, numero de escalones, distancia al puerto y temperatura media. \textbf{Analiza} cuales ayudan a un turista y cuales habria que justificar mejor.}{7}

        \begin{evidenciabox}
        \begin{itemize}
    \item Nombrar al menos 3 lugares de Almeria sobre los que se podria hacer una pregunta matematica.
    \item Distinguir entre informacion descriptiva y cuantitativa.
    \item Explicar con tus palabras para que serviria una guia turistica matematica.
\end{itemize}
        \end{evidenciabox}

        \end{document}
